\chapter{Datengrundlage}
\label{kap:daten}

Alle Parameter der folgenden Kapitel berufen sich auf Messungen am verfügbaren
Quellmaterial. Dieses Kapitel dokumentiert sie; sie sind über
\code{src/analyze\_data.py} reproduzierbar.

\section{Annotationsdefinition}
\label{sec:annotation}

Vor jeder Messung steht die Frage, was ein Objekt ist. Aus der Sichtung der
beiden handannotierten Realbilder ergibt sich eine Definition, die die gesamte
Pipeline prägt:

\begin{quote}
Annotiert ist ausschließlich die \textbf{sichtbare Stirnfläche} eines
Rundrohrprofils -- heller verzinkter Kranz, dunkle Bohrung. Rohrschäfte in
Seitenansicht, Gitterknoten und dunkle Löcher im Maschendraht sind
\textbf{keine} Objekte, obwohl sie ähnlich aussehen. Die Box umschließt den
\textbf{sichtbaren} Teil (modale Box), nicht den geschätzten vollständigen Kreis.
\end{quote}

Beide Festlegungen sind folgenreich: Die erste bestimmt, welche Strukturen als
harte Negative benötigt werden, die zweite, wie die Boxen im Generator berechnet
werden. Die Unterscheidung modal / amodal ist in der Literatur etabliert
\parencite{li2016amodal,zhu2017amodal}; amodale Boxen umschließen auch verdeckte
Teile. Erzeugte der Generator amodale Boxen, während die Realbilder modal
annotiert sind, entstünde ein systematischer Größenversatz -- ein Verstoß gegen
Anforderung~A5.

\section{Quellmaterial}
\label{sec:quellmaterial}

\absatz{Objektausschnitte.} Verfügbar sind neun aus Realaufnahmen
ausgeschnittene Rohrprofile der Größe $28\times30$ bis $43\times38$ Pixel. Sechs
zeigen eine dunkle Bohrung (offenes Rohr), drei eine helle (Sonnenlicht fällt in
die Röhre oder das Ende ist gestaucht); beide Formen kommen in den Realbildern
vor, weshalb gleichverteilt aus allen neun gezogen wird. Die aus dem
Ellipsenmodell gewonnenen Achsverhältnisse streuen zwischen 0,83 und 1,14
(Abschnitt~\ref{sec:matting}). Kritisch ist Befund~1 in Tabelle~\ref{tab:befunde}: Alle neun sind zwar RGBA-Dateien, ihr Alphakanal ist
aber durchgängig 255 -- es handelt sich um rechteckige Ausschnitte, bei denen
13--35\,\% der Fläche nicht zum Objekt gehören.

\absatz{Hintergründe.} Von den sieben Hintergrundaufnahmen sind vier
\emph{Clean Plates}: Aufnahmen derselben Objektart, aus denen die Stirnflächen
herausgeschnitten wurden -- Wände aus Rohrschäften \emph{ohne} Stirnflächen, mit
den gemessenen Rohrdurchmessern 16/18/21/20\px{}. Sie liefern gleichzeitig den
photometrisch passenden Kontext, das Rohmaterial der Okkluder-Bibliothek und
echte Negativbeispiele und erhalten deshalb sechsfaches Ziehgewicht. In der
Terminologie von Gleichung~\eqref{eq:bendavid} senken sie unmittelbar den
Divergenzterm: Der Hintergrund ist nicht \emph{ähnlich} der Zieldomäne, er
\emph{ist} sie. Die drei Fremdszenen (Kieswerk, Betonplatten, Betondeck mit
Maschendrahtzaun) sind primär als hartes Negativmaterial wertvoll -- der
Maschendraht ist genau die Struktur, die nach der Annotationsdefinition kein
Objekt ist und dennoch als Falschpositiv droht. Schwäche des Materials: Alle vier
Clean Plates zeigen dieselbe Stapelart bei bedecktem bis diffusem Licht -- kein
harter Sonnenschein, keine Nässe, kein Schnee.

\section{Bewertungsdaten und ihre Statistik}
\label{sec:valtest}

Zur Bewertung stehen zwei vollständig handannotierte Realaufnahmen zur
Verfügung: \code{val} ($2459\times1174$, 88 Objekte) für die Modellauswahl und
\code{test} ($1600\times900$, 86 Objekte) ausschließlich für die Endbewertung.

\begin{table}[htbp]
\centering
\caption[Gemessene Statistik der realen Labels]{Gemessene Statistik der realen
Labels ($n = 174$). Als Objektdurchmesser dient $\sqrt{\text{Fläche}}$; alle
Abstände sind in Objektdurchmessern (\Oe) angegeben und damit skaleninvariant
direkt als Generatorparameter verwendbar. NN = nächster Nachbar.}
\label{tab:statistik}
\footnotesize
\begin{tabular}{@{}lll@{}}
\toprule
\textbf{Größe} & \textbf{\code{val}} & \textbf{\code{test}} \\
\midrule
Durchmesser (px) & 20,6--35,8, \textbf{Median 27,4} & 12,2--31,9, \textbf{Median 19,6} \\
Seitenverhältnis $w/h$ & p10 0,97, Median 1,08, p90 1,31 & p10 0,85, Median 1,09, p90 1,22 \\
Objektfläche gesamt & 2,48\,\% des Bildes & 2,69\,\% \\
am Bildrand angeschnitten & 0\,\% & 5,8\,\% \\
\midrule
NN-Abstand (in \Oe) & p10 0,86, Median 1,09, p90 1,78 & p10 0,83, Median 1,24, p90 2,51 \\
Nachbarn innerhalb $5\,\Oe$ & Median 10, max 14 & Median 10, max 20 \\
berührt NN ($<1{,}2$\,\Oe) & 61,4\,\% & 46,5\,\% \\
Größenverhältnis zum NN & Median 1,05, p90 1,14 & Median 1,05, p90 1,17 \\
Nachbarrichtung & quasi isotrop, Häufung $90$--$150^\circ$ & quasi isotrop \\
\bottomrule
\end{tabular}
\end{table}

Bei 20--36\px{} Kantenlänge in Bildern von 1600--2459\px{} Breite handelt es sich
um Kleinobjekte im Sinne der COCO-Konvention \parencite{lin2014coco}. Nach
Gleichung~\eqref{eq:deff} schrumpfen sie bei \code{imgsz}=640 auf 5--9\px{} und
liegen damit unter der Auflösung der P3-Ebene -- die Inferenzauflösung wird
deshalb in Abschnitt~\ref{sec:imgsz} systematisch untersucht.

\section{Befunde und die daraus abgeleiteten Parameter}
\label{sec:messungmapping}

Tabelle~\ref{tab:befunde} fasst zusammen, welcher Befund welchen
Pipelineparameter festlegt. Sie ist die Belegbrücke zu
Kapitel~\ref{kap:pipeline} und macht zugleich sichtbar, welche Parameter
\emph{nicht} durch eine Messung gedeckt sind.

\begin{table}[htbp]
\centering
\caption[Befunde und abgeleitete Parameter]{Befunde der Datensichtung und die
daraus abgeleiteten Parameter. Belegart: \emph{M} = Messung an den Realdaten,
\emph{P} = physikalische Randbedingung, \emph{L} = Literatur, \emph{A} =
Annahme.}
\label{tab:befunde}
\footnotesize
\begin{tabularx}{\textwidth}{@{}L{0.30\textwidth}XL{0.05\textwidth}@{}}
\toprule
\textbf{Befund / Messwert} & \textbf{abgeleiteter Parameter} & \textbf{Art} \\
\midrule
\textbf{1} Alphakanal durchgängig 255; 13--35\,\% Fremdkontext &
eigene Matting-Stufe mit weicher Alphamaske (Signatur~S1) & P \\
\textbf{2} sechs dunkle, drei helle Bohrungen; Achsverhältnis der Ausschnitte
0,83--1,14 &
Gleichverteilung über alle neun Ausschnitte; Ovalisierung überhaupt vorsehen & A/M \\
\textbf{3} vier Clean Plates der Zieldomäne; Maschendraht und Schotter in den
Fremdszenen &
80\,\% Ziehgewicht auf Clean Plates; 10\,\% Negativbilder & M \\
\textbf{4} Rohrdurchmesser der Clean Plates 16--21\px{} &
Kopplung von Zoom und Objektgröße & P \\
\textbf{5} Objektdurchmesser 12,2--35,8\px{} &
Generatorbereich 12--44\px{} (reale Spanne plus 25\,\% Reserve) & M \\
\textbf{6} Seitenverhältnis p10--p90 = 0,85--1,31 &
Ovalisierung bei 57\,\% der Instanzen, oblat:prolat = 45:12 & M \\
\textbf{7} Größenverhältnis zum Nachbarn im Median 1,05 &
eine Basisgröße pro Cluster, je Instanz nur $\pm 8\,\%$ & M \\
\textbf{8} hinter jeder Stirnfläche führt ein Rohrschaft weg &
Rohrschaft bei 88\,\% der Instanzen (Tabelle~\ref{tab:parameter}, T5) & P \\
\textbf{9} NN-Abstand Median 1,09--1,24\,\Oe; 5,8\,\% angeschnittene Labels in
\code{test} &
Gitterabstände des Layoutmodells; Clusterzentrum darf außerhalb der Leinwand
liegen, Randregel 15\,\% & M \\
\midrule
\textbf{10} keine Messung verfügbar &
Anteil verformter Profile (18\,\%), Schrägansicht (30\,\%), Schlagschatten
(49\,\%) & A \\
\textbf{11} je \emph{ein} Bild pro Split &
Bootstrap-Konfidenzintervalle statt Punktschätzungen
(Abschnitt~\ref{sec:bootstrap}) & M \\
\bottomrule
\end{tabularx}
\end{table}

Befund~11 ist die größte methodische Schwäche der Arbeit: Je ein Bild pro Split
ist für eine belastbare Metrikangabe zu wenig, und beide Bilder zeigen dieselbe
Objektart bei ähnlichem Wetter. Alle berichteten Zahlen sind entsprechend als
Punktschätzungen mit erheblicher Unsicherheit zu lesen.
